% Laget automatisk av elevark/lag_elevark.py fra opplegg/kvantetilstander.html. Endringer her blir overskrevet.
\documentclass[11pt]{article}
\usepackage{elevark}
\begin{document}
\elevarkhode{Fysikk 2 · Kvantefysikk}{Konsekvenser av kvantetilstander – arbeidsark}{https://fysikklab.no/opplegg/kvantetilstander.html}{Opplegget på nett}

\section*{Oppgaver}

Bruk $h = 6{,}63 \cdot 10^{-34}$ Js, $m_{\mathrm{e}} = 9{,}11 \cdot 10^{-31}\,\mathrm{kg}$, $m_{\mathrm{p}} = 1{,}67 \cdot 10^{-27}\,\mathrm{kg}$, $c = 3{,}00 \cdot 10^{8}\,\mathrm{m/s}$ og 1~eV = $1{,}60 \cdot 10^{-19}\,\mathrm{J}$.

\begin{oppgaveliste}
\oppgave \fase{gjett}{Gjett}En ball ligger på bunnen av en brønn. Kan den plutselig dukke opp utenfor? Hva om ballen var et elektron, og brønnen var et atom?
\svar
\oppgave \fase{forklar}{Forklar}Vis at en partikkel i en uendelig dyp brønn med bredde $L$ har energiene $E_{n} = n^{2}h^{2}/(8mL^{2})$. Bruk at bølgen må være en stående bølge med $λ = 2L/n$.
\svar
\oppgave \fase{test}{Regn}Et elektron er fanget i en brønn på $1{,}0 \cdot 10^{-10}\,\mathrm{m}$. Regn ut de tre laveste energinivåene i eV. Hvilken bølgelengde får et foton som sendes ut når elektronet går fra $n = 2$ til $n = 1$?
\svar
\oppgave \fase{test}{Regn}Et proton er fanget i en brønn på 10 fm, omtrent som en atomkjerne. Regn ut den laveste energien i MeV. Hva slags stråling sendes ut ved overganger i kjerner?
\svar
\oppgave \fase{forklar}{Forklar}Hvorfor kan det bare være to elektroner i hver orbital? Hva ville skjedd med atomene hvis elektroner var bosoner?
\svar
\oppgave \fase{test}{Regn}En rød lysdiode sender ut lys med bølgelengde 630~nm, og en blå 470~nm. Anslå energigapet i halvlederen for hver av dem, i eV.
\svar
\oppgave \fase{forklar}{Vurder}Alice og Bob deler et par sammenfiltrede fotoner, og Alice reiser til Alpha Centauri. Når hun måler polarisasjonen til sitt foton, er polarisasjonen til Bobs foton bestemt. Kan de bruke dette til å sende meldinger fortere enn lyset?
\svar
\oppgave \fase{utfordring}{Utfordring}Temperaturen i sentrum av sola er omtrent $1{,}5 \cdot 10^{7}$ $K$. Regn ut den typiske termiske energien kT $(k = 1{,}38 \cdot 10^{-23} J/K)$ i keV. Den elektriske barrieren mellom to protoner er i størrelsesorden 1~MeV. Hvordan kan fusjon likevel skje?
\svar
\end{oppgaveliste}

\end{document}
